\subsubsection{Định lý biểu diễn}

Lưu ý rằng, bỏ qua hệ số dịch chuyển $b$, nghiệm của SVM
có thể được biểu diễn dưới dạng tổ hợp tuyến tính của các hàm
$K(x_i,\cdot)$, trong đó $x_i$ là các điểm dữ liệu huấn luyện.
Định lý sau đây, được gọi là \emph{định lý biểu diễn}
(\emph{Representer Theorem}), cho thấy rằng đây là một tính chất
tổng quát đúng cho nhiều lớp bài toán tối ưu,
bao gồm cả SVM không có hệ số dịch chuyển.

\begin{thm}[Định lý biểu diễn]
\label{thm:representer}

Cho
$
K:\mathcal{X}\times\mathcal{X}\to\mathbb{R}
$
là một hàm nhân đối xứng xác định dương (PDS),
và $\mathbb{H}$ là không gian Hilbert nhân tái tạo
(RKHS) tương ứng với $K$.

Khi đó, với mọi hàm không giảm
$
G:\mathbb{R}\to\mathbb{R}
$
và mọi hàm mất mát
$
L:\mathbb{R}^{m}\to\mathbb{R}\cup\{+\infty\},
$
bài toán tối ưu
\begin{equation}
\argmin_{h\in\mathbb{H}} F(h)
=
\argmin_{h\in\mathbb{H}}
G\bigl(\|h\|_{\mathbb{H}}\bigr)
+
L\bigl(h(x_1),\dots,h(x_m)\bigr)
\end{equation}

luôn tồn tại một nghiệm có dạng
\begin{equation}
h^{*}
=
\sum_{i=1}^{m}
\alpha_i K(x_i,\cdot).
\end{equation}

Nếu $G$ còn được giả sử là hàm tăng,
thì mọi nghiệm của bài toán tối ưu
đều có dạng trên.

\end{thm}

\begin{proof}
Gọi
$\mathbb{H}_1
=
\mathrm{span}\{K(x_i,\cdot): i \in [m]\}$.
Khi đó mọi $h \in \mathbb{H}$ có phân rã trực giao
$
h = h_1 + h^\perp,
$
với
$
\mathbb{H}
=
\mathbb{H}_1 \oplus \mathbb{H}_1^\perp.
$
trong đó $\oplus$ là tổng trực tiếp.
Vì $G$ là hàm không giảm nên
\[
G(\|h_1\|_{\mathbb{H}})
\le
G\!\left(
\sqrt{
\|h_1\|_{\mathbb{H}}^2
+
\|h^\perp\|_{\mathbb{H}}^2
}
\right)
=
G(\|h\|_{\mathbb{H}}).
\]

Theo tính chất tái sinh, với mọi $i \in [m]$,
\[
h(x_i)
=
\langle h, K(x_i,\cdot)\rangle
=
\langle h_1, K(x_i,\cdot)\rangle
=
h_1(x_i).
\]

Do đó,
\[
L(h(x_1),\ldots,h(x_m))
=
L(h_1(x_1),\ldots,h_1(x_m)),
\]
suy ra
$
F(h_1) \le F(h).
$

Điều này chứng minh phần đầu của định lý.

Nếu $G$ còn là hàm tăng nghiêm ngặt và
$
\|h^\perp\|_{\mathbb{H}} > 0,
$
thì
$
F(h_1) < F(h),
$
do đó mọi nghiệm tối ưu đều phải thuộc
$\mathbb{H}_1$.
\end{proof}